% ============================================================
% ANEXO B - Respuestas preparadas
% Bloques: Validacion y Metricas (Maria Escudero) + Especificacion y Gestion (Danela Arteaga)
% Fragmento para \input en el informe final
% ============================================================
\subsection*{Fundamentos}

\noindent\textbf{1. ¿Dónde está la frontera del sistema y qué quedó fuera de forma explícita?}\\
La frontera es la aplicación web del SGA: gestión de parcelas y cultivos, registro agrícola (cosecha, labor, enfermedad), inventario y compras con alerta de stock mínimo, planificación de tareas, reportes e ingresos, y los dos componentes de IA especificados en esta práctica (IA-01 recomendador y IA-02 detector por imagen). Quedó fuera de forma explícita: la app móvil nativa y cualquier sensor o estación meteorológica (no tienen requisito asociado en el catálogo), el módulo contable completo ---el control financiero lo lleva el contador externo; solo entra el dato operativo que lo alimenta---, pasarelas de pago y logística de venta post-cosecha, y el entrenamiento y despliegue productivo de los modelos: aquí especificamos requisitos, fichas de datos y plan de monitoreo de la IA, no su implementación.

\medskip
\noindent\textbf{2. ¿Qué muestra el diagrama de contexto que no muestra el límite del sistema?}\\
El diagrama de contexto responde quién interactúa con el sistema y qué información cruza la frontera: el administrador y el trabajador agrícola como actores operativos, el contador externo que recibe el reporte de pérdidas, los proveedores de insumos y las partes interesadas institucionales (equipo y docente). El límite, en cambio, define qué procesa el software por dentro. La diferencia se ve con un ejemplo: RF-19 produce el reporte de pérdidas económicas por plagas como salida hacia el contador, pero todo el proceso contable posterior queda fuera; el contexto modela ese intercambio de información, mientras que el límite declara que no es responsabilidad del SGA procesarlo.

\medskip
\noindent\textbf{3. ¿Qué modelo de procesos siguieron y por qué ese y no otro?}\\
Un ciclo incremental alineado con los procesos técnicos de ISO/IEC/IEEE 15288 y 12207: definición de requisitos de interesados (entrevistas y modelos i* en la primera entrega) y análisis de requisitos del sistema (ERS línea base y su corrección), cerrando cada iteración una línea base versionada en Git. Lo elegimos porque el alcance creció durante el propio proyecto ---los componentes de IA entraron recién en el Paso~3 de la PE5--- y el incremental permitió validar temprano con el administrador, corregir con evidencia versionada (nueve observaciones de línea base y cinco hallazgos de re-inspección documentados) y re-auditar las métricas sobre cada versión sin reescribir el documento completo.

\subsection*{Elicitación}

\noindent\textbf{4. ¿Qué técnica les aportó requisitos que ninguna otra habría descubierto?}\\
La observación directa en campo. Las bitácoras y registros fotográficos de las jornadas mostraron que el control de maleza a machete era la labor que más tiempo consumía y que las instrucciones se repartían de forma verbal, sin rastro verificable. En la encuesta digital el trabajador había respondido que ``recibe órdenes'', pero solo al observar la jornada quedó claro el problema de fondo: no existía trazabilidad de qué tarea se asignó, a quién y si se cumplió. De ahí nacieron RF-08 a RF-10 (planificación y asignación de tareas con evidencia de cumplimiento) como prioridad Debe tener; ninguna otra técnica hubiera dado esa priorización con fundamento.

\medskip
\noindent\textbf{5. ¿Qué stakeholder quedó subrepresentado y cómo compensaron su ausencia?}\\
El contador externo: gestiona las finanzas de la finca, pero no fue entrevistado directamente; su necesidad nos llegó mediada por el administrador en la entrevista principal. Compensamos su ausencia de dos maneras. Primero, delimitando el alcance para no decidir por él: el módulo contable completo quedó fuera del sistema. Segundo, modelándolo como actor externo con una única interacción bien definida, recibir el reporte de pérdidas por plagas (RF-19), que es precisamente el dato que el administrador señaló como el más valioso (EV-01 [08:17]). Registramos la lección: un interesado no entrevistado debe entrar al modelo con el menor compromiso contractual posible, hasta poder validar con él.

\medskip
\noindent\textbf{6. Muéstrenme un requisito y díganme de qué evidencia nació.}\\
RF-19, reporte de pérdidas económicas por plagas y enfermedades, nace textualmente de EV-01 [08:17]: el administrador declaró que saber cuánto pierde por plaga es la información que más le cuesta obtener hoy y la que más usaría para decidir inversiones fitosanitarias. La ficha del requisito cita esa evidencia en su campo Origen, la matriz de trazabilidad lo enlaza a la misma fuente, y su criterio de verificación ---totales por cultivo y período contrastables contra los registros de cosecha e incidencias--- se deriva del enunciado literal del entrevistado, no de una suposición nuestra.

\subsection*{Validación}

\noindent\textbf{10. ¿Cuántos defectos halló la inspección y cuáles siguen abiertos?}\\
Entre el ciclo de revisión de la Entrega 2 y la re-inspección interna registramos nueve observaciones corregidas entre entregas (LB-01 a LB-09, en \texttt{registro\_defectos\_lineabase\_1B\_2A.csv}) y cinco hallazgos únicos en la re-inspección de esta práctica (DR-01 a DR-05, en \texttt{registro\_reinspeccion\_PE5.csv}). De estos últimos, los cinco quedaron cerrados al final de la práctica: DR-01 (crítico de trazabilidad) se cerró con las 20 filas nuevas de la matriz y las cuatro historias de usuario faltantes; los otros cuatro quedaron cerrados antes: tres corregidos sobre el texto (RF-25, RNF-05, RNF-15) y uno justificado documentalmente (TR-19, requisito Could Have).

\medskip
\noindent\textbf{11. ¿Qué defecto crítico se les escapó y cómo lo detectaron después?}\\
El más grave que se nos pasó fue DR-01: añadimos once requisitos funcionales en la Entrega 3 con sus fichas completas, pero nadie verificó que la matriz de trazabilidad creciera junto con el catálogo; la matriz quedó corta en la fila TR-22. Lo detectamos en la PE5 al recorrer la matriz fila a fila contra la plantilla de diez columnas exigida, comparando los identificadores presentes contra el catálogo completo de 42 requisitos: esa comparación directa es la que hizo visible el faltante. La lección que dejamos registrada es que cada alta de requisito debe arrastrar su fila de matriz en el mismo commit, no en una pasada posterior.

\medskip
\noindent\textbf{12. ¿Cómo saben que las correcciones no introdujeron nuevos defectos?}\\
Porque la re-inspección incluyó una pasada completa posterior sobre lo corregido. Primero, las cuatro correcciones de texto son acotadas y verificables por diff: el criterio reescrito de RF-25, la referencia actualizada de RNF-05 y la enumeración de métricas de RNF-15 están en commits revisables del repositorio. Segundo, recalculamos las dos métricas afectadas sobre la versión corregida (M5 bajó de 3,50 a 3,00 y M6 de 0,119 a 0,024) usando los conteos publicados antes de calcular en \texttt{conteos\_base\_auditoria.csv}, así que cualquier regresión habría aparecido como inconsistencia entre conteo y texto al re-auditar. Tercero, el cierre de DR-01 (matriz ampliada a 62 filas) se hizo sin tocar ninguna de las 42 filas originales de contenido, solo se les agregaron columnas nuevas; un diff del archivo lo confirma.

\subsection*{Métricas}

\noindent\textbf{19. ¿Qué métrica salió peor y qué hicieron al respecto?}\\
La primera medición dejó dos métricas fuera de referencia. M6 (corrección) dio 0,119 defectos residuales por requisito contra un máximo de 0,05: aplicamos las correcciones DR-02 a DR-04, justificamos DR-05 y quedó 0,024 con un único residual gestionado. M5 (modificabilidad) dio 3,50 requisitos afectados por cambio contra un máximo de 3,0: reespecificamos RF-12, RF-18 y RF-27 como vistas derivadas de la consulta por período RF-07 para quitarles dependencia directa de RF-05, y la re-medición sobre la misma muestra dio 3,00. La que quedó peor al inicio fue M4 (trazabilidad), en 66,7\,\% frente a un mínimo de 90\,\%; su cierre era estructural y quedó completado en el Paso 2 (ver bloque Gestión más abajo): la matriz pasó a cubrir el 100\,\% de los 42 requisitos del ERS.

\medskip
\noindent\textbf{20. Enséñenme los conteos con los que calcularon la verificabilidad.}\\
Están publicados en \texttt{09\_Cierre\_PE5/conteos\_base\_auditoria.csv} y coinciden con el Anexo A del informe: 42 requisitos en total (27 funcionales RF-01 a RF-27 y 15 no funcionales RNF-01 a RNF-15), de los cuales los 42 tienen criterio de verificación con métrica, unidad, valor objetivo y método de medición tras la corrección LB-01, lo que da $42/42 = 100\%$ contra una referencia de $\geq 90\%$. De esos 42, los dieciséis requisitos de prioridad Debe tener tienen además criterios de aceptación en formato Gherkin en la Sección 3.4 del ERS y en \texttt{04\_Trazabilidad/HU\_criterios\_aceptacion.md}.

\subsection*{Especificación}

\noindent\textbf{7. Elijan un RNF y expliquen cómo se comprueba objetivamente.}\\
Tomamos RNF-09 (Tolerancia a interrupciones durante el registro). Su verificación es objetiva porque no depende de una opinión: se simula una interrupción de conexión durante el envío de un formulario de registro (cosecha, labor o enfermedad) en 20 intentos, y se cuenta el porcentaje de esos intentos en que el sistema conserva sin pérdida los datos ya ingresados y permite reintentar. El valor objetivo es 100\,\%. Dos personas evaluadoras que ejecuten la misma prueba sobre los mismos 20 intentos deben llegar al mismo veredicto, porque el criterio no exige interpretación: los datos están o no están tras la interrupción.

\medskip
\noindent\textbf{8. ¿Qué requisito les costó más redactar de forma verificable?}\\
RF-26 (detección temprana de enfermedad mediante IA). A diferencia de un requisito transaccional como RF-01, aquí el resultado es probabilístico: no se puede exigir 100\,\% de acierto porque el propio administrador reconoce que la inteligencia artificial tiene margen de error (EV-01 [09:10]). Tuvimos que fijar un umbral realista (exactitud mínima del 70\,\% sobre un conjunto de validación de al menos 30 casos etiquetados) y, sobre todo, enlazarlo con RNF-15 (explicabilidad): sin la explicación causal simultánea a la alerta, el requisito sería verificable en el papel pero inútil en la práctica, porque el administrador condiciona su confianza a poder verificar la alerta antes de actuar.

\medskip
\noindent\textbf{9. ¿Por qué este caso de uso incluye a aquel y no lo extiende?}\\
En el diagrama general (UC-01), UC-15 (Consultar recomendaciones de IA) incluye \emph{Generate Crop Recommendations}, que a su vez incluye \emph{Analyze Crop Data}: es una relación \texttt{include} porque analizar los datos del cultivo es un paso obligatorio y siempre presente para producir una recomendación, no un comportamiento opcional. En cambio, UC-16 (Registrar enfermedad del cultivo) se relaciona con \emph{Generate Alert} mediante \texttt{extend}: generar una alerta solo ocurre condicionalmente, cuando la enfermedad registrada tiene gravedad alta; en el resto de los casos, el flujo de UC-16 se completa sin ese paso adicional. La diferencia es esa: \texttt{include} para un paso siempre necesario, \texttt{extend} para un comportamiento condicional que se activa solo bajo una regla de negocio específica.

\subsection*{Gestión}

\noindent\textbf{13. Tomen un RF y digan qué se rompería si cambia mañana.}\\
Tomamos RF-06 (gestión de inventario). Según la matriz de dependencias usada para calcular M5 (\texttt{09\_Cierre\_PE5/calculo\_M5\_modificabilidad.csv}), RF-06 es entrada directa de otros tres requisitos: RF-21 (planificación de compra, que selecciona el insumo del inventario), RF-22 (entrega de insumo, que descuenta cantidad del inventario) y RF-25 (alerta por disponibilidad crítica, que se dispara cuando un insumo cae bajo su umbral mínimo). Cambiar la estructura de datos de \texttt{InventoryItem} —por ejemplo, el campo de umbral mínimo— obligaría a revisar los tres RF dependientes, no solo el propio RF-06.

\medskip
\noindent\textbf{14. ¿Qué RFC rechazó el CCB?}\\
El registro de defectos de línea base (\texttt{registro\_defectos\_lineabase\_1B\_2A.csv}) documenta cuatro observaciones corregidas entre la Entrega 2 y la Entrega 3 (LB-01 a LB-04), todas aceptadas y aplicadas; no consta ningún RFC rechazado por el Comité de Control de Cambios en ese ciclo. Si el docente pregunta por una propuesta específicamente rechazada, hay que confirmarlo con María, que es quien recuperó y versionó ese registro.

\medskip
\noindent\textbf{15. ¿Cuál es la línea base vigente y dónde está etiquetada?}\\
La línea base vigente al momento de esta práctica es \texttt{ERS\_SRS\_2A\_v1.0.tex}, corregida el 22-08-2026 con los cinco hallazgos de la re-inspección (DR-01 a DR-05) y ampliada con las 20 filas nuevas de trazabilidad. La versión de cierre de la PE5 será \texttt{ERS\_SRS\_2B\_v2.0}, que consolida Robinson; su etiqueta de Git y el commit exacto deben confirmarse el día de la defensa, una vez subida esa versión final.
